% !TEX program = xelatex
% Compile with: xelatex report.tex
\documentclass[12pt,a4paper]{article}

\usepackage{fontspec}
\usepackage{polyglossia}
\setmainlanguage{ukrainian}
\setotherlanguage{english}

% Use system fonts that support Cyrillic. Times New Roman and Courier New
% both ship with macOS and contain full Cyrillic coverage.
\setmainfont{Times New Roman}
\setmonofont[Scale=0.88]{Courier New}
\newfontfamily\cyrillicfonttt[Scale=0.88]{Courier New}

\usepackage[margin=2.2cm]{geometry}
\usepackage{graphicx}
\usepackage{hyperref}
\usepackage{xcolor}
\usepackage{listings}
\usepackage{caption}
\usepackage{float}

\hypersetup{
  colorlinks=true,
  linkcolor=black,
  urlcolor=blue,
  citecolor=blue
}

\definecolor{lsbg}{gray}{0.97}
\definecolor{lskw}{RGB}{30,80,160}
\definecolor{lscmt}{RGB}{90,130,90}
\definecolor{lsstr}{RGB}{150,60,30}

\lstdefinelanguage{arm64asm}{
  morekeywords={stp,ldp,str,ldr,strb,ldrb,ldrh,mov,movk,movz,add,sub,subs,mul,madd,msub,udiv,
                lsl,lsr,asr,and,orr,eor,cmp,cmn,ccmp,cbz,cbnz,b,bl,ret,
                adrp,adr,csel,cset,tst,b.eq,b.ne,b.lo,b.hi,b.hs,b.ls,b.lt,b.gt,b.ge,b.le},
  morekeywords=[2]{.text,.data,.bss,.section,.align,.equ,.globl,.extern,.byte,.asciz,.space},
  morecomment=[l]{//},
  morecomment=[s]{/*}{*/},
  morestring=[b]",
  sensitive=true,
}

\lstset{
  basicstyle=\ttfamily\footnotesize,
  backgroundcolor=\color{lsbg},
  keywordstyle=\color{lskw}\bfseries,
  keywordstyle=[2]\color{lskw},
  commentstyle=\color{lscmt}\itshape,
  stringstyle=\color{lsstr},
  breaklines=true,
  numbers=left,
  numberstyle=\tiny\color{gray},
  frame=single,
  rulecolor=\color{gray!40},
  framexleftmargin=2mm,
  xleftmargin=5mm,
  showstringspaces=false,
  tabsize=4,
  captionpos=b
}

% Placeholder command for screenshots. Replace path/to/... after dropping
% an image into report/images/ with the given file name.
\newcommand{\screenshot}[3][0.9\linewidth]{%
  \begin{figure}[H]
    \centering
    \IfFileExists{images/#2}{%
      \includegraphics[width=#1]{images/#2}%
    }{%
      \fbox{\parbox[c][5cm][c]{#1}{\centering\itshape
        Сюди помістіть скрін \texttt{images/#2}}}%
    }%
    \caption{#3}
    \label{fig:#2}
  \end{figure}
}

\begin{document}

% =====================================================================
% Title page
% =====================================================================
\begin{titlepage}
  \centering
  \vspace*{1.5cm}
  {\Large Міністерство освіти і науки України\par}
  \vspace{0.3cm}
  {\Large Національний технічний університет України\par}
  {\Large «Київський політехнічний інститут імені Ігоря Сікорського»\par}

  \vspace{2cm}
  {\large Факультет <<Назва факультету>>\par}
  {\large Кафедра <<Назва кафедри>>\par}

  \vfill
  \rule{\linewidth}{0.8pt}\par
  \vspace{0.5cm}
  {\LARGE \bfseries Лабораторний проєкт~№2\par}
  \vspace{0.3cm}
  {\large <<Робота з файловою системою>>\par}
  \vspace{0.3cm}
  {\large Варіант~4: перегляд PNG, лічильник JPG/PNG у папці,\\
          ім'я та дата створення для інших файлів\par}
  \vspace{0.5cm}
  \rule{\linewidth}{0.8pt}\par
  \vfill

  \begin{flushright}\large
    Виконав(ла): студент(ка) групи \textbf{<<група>>}\\
    \textbf{<<Прізвище Ім'я По батькові>>}\\[0.8em]
    Перевірив(ла): \textbf{<<Посада, ПІБ>>}
  \end{flushright}

  \vfill
  {\large Київ~--- \the\year}
\end{titlepage}

\tableofcontents
\newpage

% =====================================================================
\section{Умова задачі}
% =====================================================================

Необхідно написати програму \emph{Переглядач файлової системи} (ПФС)
із двома панелями:

\begin{itemize}
  \item \textbf{Ліва панель} --- список файлів і папок активної директорії.
  \item \textbf{Права панель} --- інформація відповідно до варіанту.
\end{itemize}

ПФС має дозволяти підніматися в надпапку і заходити в обрану папку.
Реалізація --- на розсуд студента.

\paragraph{Варіант~4.} У правій панелі:
\begin{itemize}
  \item для \texttt{.png}-файлів --- перегляд зображення;
  \item для папок --- кількість \texttt{.jpg}/\texttt{.png}-файлів усередині;
  \item для інших файлів --- ім'я та дата створення.
\end{itemize}

У цьому проєкті обсяг відображуваної метаінформації розширений:
для будь-якого файлу додатково виводяться права доступу,
кількість hard-links, власник, inode, а також три часові мітки
(\texttt{atime}, \texttt{mtime}, \texttt{btime}).

% =====================================================================
\section{Середовище та технології}
% =====================================================================

\begin{itemize}
  \item Платформа: macOS 15+, архітектура Apple~Silicon (\texttt{arm64}).
  \item Ядро проєкту реалізовано мовою асемблера \textbf{ARM64 (AAPCS64)}.
  \item Графічний інтерфейс --- \textbf{Objective-C~+~AppKit} (Cocoa),
        компіляція з \texttt{-fobjc-arc}.
  \item Збірка --- \texttt{clang}, керування через \texttt{make}.
\end{itemize}

% =====================================================================
\section{Архітектура}
% =====================================================================

Проєкт розділено на два шари:

\begin{description}
  \item[Низькорівневий asm-шар] реалізує все, що стосується файлової
        системи, розбору PNG та форматування:
        \begin{itemize}
          \item \texttt{src/fs.s} --- перелік директорії (\texttt{fs\_list\_dir}),
                підйом у надпапку (\texttt{path\_parent\_inplace}),
                підрахунок jpg/png файлів (\texttt{fs\_count\_jpg\_png}),
                базові метадані файлу (\texttt{pfs\_file\_size},
                \texttt{pfs\_file\_mtime\_iso}), розгорнутий вивід
                (\texttt{pfs\_format\_fullmeta}).
          \item \texttt{src/path\_join.s} --- склеювання шляху.
          \item \texttt{src/png\_ihdr.s} --- зчитування, валідація
                PNG-сигнатури й розбір IHDR-чанка.
          \item \texttt{src/pfs\_fmt.s} --- дрібні хелпери форматування:
                \texttt{pfs\_str\_append}, \texttt{pfs\_u32\_append},
                \texttt{pfs\_u64\_append}.
        \end{itemize}
  \item[GUI-шар на Objective-C] відповідає лише за AppKit-подання
        (\texttt{NSWindow}, \texttt{NSTableView}, \texttt{NSImageView},
        \texttt{NSTextView}, кнопки, обробники):
        \begin{itemize}
          \item \texttt{src/main.m} --- створення \texttt{NSApplication}
                та утримання делегата.
          \item \texttt{src/AppDelegate.m} --- побудова вікна,
                двох колонок, обробка подій, виклики asm-функцій.
        \end{itemize}
\end{description}

\subsection{Взаємодія між шарами}

Декларації низькорівневих функцій зібрано в \texttt{src/pfs\_core.h},
що його імпортує ObjC. Дані про вміст папки лежать у BSS-сегменті
(\texttt{g\_entries}, \texttt{g\_entry\_count}), описаних у
\texttt{fs.s}. GUI перечитує цей масив після кожного
\texttt{fs\_list\_dir} і використовує його як джерело для
\texttt{NSTableViewDataSource}.

\subsection{Використані системні виклики}

\begin{itemize}
  \item \texttt{opendir} / \texttt{readdir} / \texttt{closedir} ---
        обхід директорій.
  \item \texttt{lstat} --- метадані файлу.
  \item \texttt{open} / \texttt{read} / \texttt{close} ---
        читання PNG-сигнатури й IHDR-чанка.
  \item \texttt{localtime\_r}, \texttt{strftime} ---
        форматування часу у форматі ISO\,8601.
  \item \texttt{qsort} --- упорядкування записів директорії.
\end{itemize}

% =====================================================================
\section{ARM64 ABI (AAPCS64) і характерні пастки}
% =====================================================================

У процесі реалізації виявились класичні проблеми з дотриманням
ABI AAPCS64, що не помітні у коротких демонстраціях, але з'являються
у великих ObjC-програмах, де компілятор активно використовує
callee-saved регістри для збереження \texttt{self}:

\begin{itemize}
  \item Регістри \texttt{x19}--\texttt{x28} --- callee-saved:
        якщо асемблерна функція їх перезаписує, вона \textbf{зобов'язана}
        зберегти/відновити на стеку.
  \item Під ARC ObjC-компілятор розміщує \texttt{self} у \texttt{x19}
        на увесь час виконання метода. Підсумок: асемблерна функція,
        що написала в \texttt{x19}, ламає \texttt{self} у caller-і.
  \item На Darwin~ARM64 варіадичні аргументи передаються через стек,
        а не через регістри --- це зумовило перенос
        \texttt{png\_format\_ihdr\_info} з використання \texttt{snprintf}
        на власні хелпери форматування без варіадів.
\end{itemize}

Тому всі asm-функції у цьому проєкті або не торкаються
\texttt{x19}--\texttt{x28} зовсім, або зберігають їх повністю
(\texttt{stp/ldp} парами).

% =====================================================================
\section{Ключові фрагменти коду}
% =====================================================================

\subsection{Склеювання шляху (\texttt{path\_join.s})}

Функція виконується без жодного зовнішнього виклику; працює в
caller-saved регістрах, що дозволяє обійтись без кадру стеку.

\lstinputlisting[language=arm64asm,caption={\texttt{src/path\_join.s}},
                 label={lst:path_join}]
{../src/path_join.s}

\subsection{Розбір PNG (\texttt{png\_ihdr.s})}

Функція читає файл у стековий буфер (8\,КіБ), валідує сигнатуру й
структуру IHDR, витягує ширину, висоту та колірні параметри,
а далі формує текстовий опис через власні append-хелпери.

\lstinputlisting[language=arm64asm,caption={\texttt{src/png\_ihdr.s}},
                 label={lst:png_ihdr}]
{../src/png_ihdr.s}

\subsection{Форматування (\texttt{pfs\_fmt.s})}

Три універсальні хелпери для безпечного дозапису у буфер з
контролем меж; цілі числа виводяться без \texttt{snprintf},
цифра за цифрою.

\lstinputlisting[language=arm64asm,caption={\texttt{src/pfs\_fmt.s}},
                 label={lst:pfs_fmt}]
{../src/pfs_fmt.s}

\subsection{Файлова система та метадані (\texttt{fs.s})}

Основний asm-модуль: сканування директорії, підйом у надпапку,
підрахунок jpg/png, одержання розміру файлу, формування
ISO-часу та розгорнутого блоку метаданих (права, власник,
inode, три часові мітки).

\lstinputlisting[language=arm64asm,caption={\texttt{src/fs.s}},
                 label={lst:fs}]
{../src/fs.s}

\subsection{Інтеграція в GUI (\texttt{AppDelegate.m})}

ObjC-шар тільки викликає asm-функції та подає їхній вивід у
\texttt{NSTextView} / \texttt{NSImageView}.

\lstinputlisting[language={[Objective]C},
                 caption={\texttt{src/AppDelegate.m}},
                 label={lst:appdelegate}]
{../src/AppDelegate.m}

% =====================================================================
\section{Складання та запуск}
% =====================================================================

\begin{lstlisting}[language=bash,caption={Збирання}]
make        # src/*.s + src/*.m -> ./pfs
./pfs       # запуск
make clean  # прибирання
\end{lstlisting}

Makefile виконує:
\begin{itemize}
  \item \texttt{clang -arch arm64 -c} кожного \texttt{.s} файлу;
  \item \texttt{clang -arch arm64 -fobjc-arc -c} кожного \texttt{.m} файлу;
  \item лінкування з \texttt{-framework Cocoa -framework Foundation}.
\end{itemize}

% =====================================================================
\section{Скріни роботи програми}
% =====================================================================

\screenshot{main-window.png}{Стартове вікно, відкритий робочий каталог.
Ліва панель містить список, справа --- метадані.}

\screenshot{png-preview.png}{Для \texttt{.png}-файлу показано IHDR-інформацію
та самé зображення.}

\screenshot{jpeg-meta.png}{Для \texttt{.jpg}/\texttt{.jpeg} виводяться
метадані (розмір, права, inode, часи) і прев'ю.}

\screenshot{folder-count.png}{Для папки виводиться кількість
\texttt{.jpg}-/\texttt{.png}-файлів усередині.}

\screenshot{parent-navigation.png}{Навігація <<Parent folder>> --- підйом
у надкаталог.}

% =====================================================================
\section{Висновки}
% =====================================================================

У ході виконання роботи:
\begin{itemize}
  \item Реалізовано ПФС згідно з варіантом~4: перегляд \texttt{.png},
        лічильник jpg/png у папці, метадані для решти файлів.
  \item Максимум функціональності винесено на ARM64-асемблер:
        читання директорії, розбір PNG, форматування чисел і часу,
        повний блок метаданих.
  \item Objective-C використовується лише як тонкий шар для AppKit
        і не містить бізнес-логіки.
  \item Особливу увагу приділено дотриманню AAPCS64 у всіх
        асемблерних функціях --- це критично для стабільної роботи
        з ObjC ARC-кодом, що тримає \texttt{self} у \texttt{x19}.
\end{itemize}

\end{document}
